\subsubsection{Path space and invariant measure between sheets}
\label{sec:medida-retorno-hojas}

The toroidal topology of APP fixes the support of the cardinal translations
and their winding classes. The correspondence between sheets arises from
another operation on that same support: the temporal selection of the
additive or multiplicative evaluation. Its geometric typing is
\[
 \eta_{\rm av}(\Sigma)=\mathscr H_{\rm ar},\qquad
 \eta_{\rm av}(\Pi)=\mathscr H_{\rm vol}.
\]
The first sheet corresponds to the additive boundary reading; the second,
to the multiplicative interior reading. This map transports the mode labels
and preserves the incidence \eqref{eq:fav}. The square of the self-scaling
ratio is determined through the paths of that incidence.

\Needspace{8\baselineskip}
\begin{definition}[One-memory symbolic envelope]
Let
\[
 X_{\rm av}=\{(s_k)_{k\in\mathbb Z}\in
 \{\mathscr H_{\rm ar},\mathscr H_{\rm vol}\}^{\mathbb Z}:
 (F_{\rm av})_{s_k,s_{k+1}}=1\text{ for every }k\}.
\]
The index shift acts on this space. It is the smallest one-memory symbolic
system that contains the modal sequences and preserves their consecutive
pairs: it admits the two cross transitions and volumetric self-retention.
The phase, cardinal orientation, and path records remain in the TPK state
from which this envelope is obtained.
\end{definition}

\begin{proposition}[Weighting of paths between sheets]
In areal--volumetric order, the invariant measure of maximal entropy
of \(X_{\rm av}\) has transition matrix and stationary distribution
\begin{equation}
 P_{\rm av}=\begin{pmatrix}0&1\\\varphi^{-2}&\varphi^{-1}\end{pmatrix},
 \qquad
 (\mu_{\rm ar},\mu_{\rm vol})
 =\frac{(1,\varphi^2)}{1+\varphi^2}.
 \label{eq:medida-hojas-exacta}
\end{equation}
In particular, \(\mu_{\rm ar}/\mu_{\rm vol}=\varphi^{-2}\).
\end{proposition}
\begin{proof}
The incidence already constructed has positive eigenvector
\(r=(1,\varphi)^{\mathsf T}\) and eigenvalue \(\varphi\). The
normalization
\[
 (P_{\rm av})_{ij}
 =\frac{(F_{\rm av})_{ij}r_j}{\varphi r_i}
\]
gives the stated matrix. Its rows sum to one because
\(\varphi^{-2}+\varphi^{-1}=1\). The symmetry of \(F_{\rm av}\)
implies that the weights proportional to \(r_i^2\) satisfy detailed
balance; hence \(\mu P_{\rm av}=\mu\). Irreducibility and volumetric
self-retention determine a unique stationary distribution.

To verify maximality, every allowed transition satisfies
\(\log(P_{\rm av})_{ij}=\log r_j-\log r_i-\log\varphi\).
Under any invariant measure, the endpoint terms cancel on average.
If \(\nu_i\) are its vertex weights and
\(q_i(j)=\nu(s_1=j\mid s_0=i)\), for \(\nu_i>0\), then
\[
 h_\nu\leq H_\nu(s_1\mid s_0)
 =\log\varphi-\sum_{i:\nu_i>0}\nu_iD(q_i\Vert P_i)
 \leq\log\varphi,
\]
where \(D(q\Vert p)=\sum_jq_j\log(q_j/p_j)\) is the relative entropy
on the allowed pairs, with the convention \(0\log0=0\).
The Markov measure of \eqref{eq:medida-hojas-exacta} attains the bound.
Equality in the first inequality requires the Markov property;
equality in the second requires \(q_i=P_i\). Stationarity and
irreducibility then fix \(\mu\), thereby establishing uniqueness.
This is the Parry measure of the defined envelope.
\end{proof}

The frequency \((1/3,2/3)\) counts the instants of the primitive calendar;
\(\mu\) weights the paths of the envelope. Its quadratic ratio
expresses the product of the incoming and outgoing incidence amplitudes.
The memory and phase constraints of the TPK preserve their own domain;
the entropy of the envelope corresponds to the symbolic system
specified in this definition.

\subsubsection{Closure operator and limit of the marked returns}
\label{sec:operador-retorno-marcado}

The closure coordinate \(\pi\), generated by the regional reader,
enters after the construction of the modal incidence. Let
\(e_{\rm ar}=(1,0)^{\mathsf T}\),
\(e_{\rm vol}=(0,1)^{\mathsf T}\) and
\(P_i=e_ie_i^{\mathsf T}\) be the sheet projectors. The conditions
of the reader are preservation of both sheets, unit volumetric
normalization, and an areal mark equal to the closure coordinate. These
conditions determine the positive operator
\begin{equation}
 D_\pi=\pi P_{\rm ar}+P_{\rm vol}
       =\begin{pmatrix}\pi&0\\0&1\end{pmatrix}.
 \label{eq:operador-marca-pi}
\end{equation}
Diagonality expresses sheet preservation, and the two entries express
their normalizations; the uniqueness of the operator for the fixed reader
is thereby made explicit.

\Needspace{8\baselineskip}
\begin{proposition}[Ratio of boundary returns]
For \(n\geq1\), the marked evaluation
\begin{equation}
 \Theta_n=
 \frac{\langle e_{\rm ar},D_\pi F_{\rm av}^{\,n}e_{\rm ar}\rangle}
 {\langle e_{\rm vol},F_{\rm av}^{\,n}e_{\rm vol}\rangle}
 =\pi\frac{F_{n-1}}{F_{n+1}}
 \label{eq:retorno-marcado-finito}
\end{equation}
satisfies
\begin{equation}
 \lim_{n\to\infty}\Theta_n
 =\pi\frac{\mu_{\rm ar}}{\mu_{\rm vol}}
 =\frac{\pi}{\varphi^2}.
 \label{eq:retorno-marcado-limite}
\end{equation}
\end{proposition}
\begin{proof}
The diagonal entries of \(F_{\rm av}^{\,n}\) count the paths
of length \(n\) that return to their initial sheet. The power formula
already proved yields \(F_{n-1}\) and \(F_{n+1}\). The operator
\(D_\pi\) multiplies the first count by \(\pi\). Finally,
\(F_{n-1}/F_{n+1}\) is the product of two consecutive quotients
that converge to \(\varphi^{-1}\); its limit is
\(\varphi^{-2}\). The identity with the ratio of weights follows from
\eqref{eq:medida-hojas-exacta}.
\end{proof}

The closure operator appears only once, in the terminal evaluation
of the return. The transfer \(F_{\rm av}\) remains determined by
the calendar. The information from two readers of the same state is
thus composed: the modal incidence determines the quadratic weighting,
and the closure region contributes its coordinate at the boundary.

The composition also preserves refinement. For positive cylinders
\(I=[a,b]\) and \(J=[c,d]\), the image of the reader
\(\mathcal R(x,y)=x/y^2\) is the interval
\[
 \mathcal Q(I,J)=\left[\frac{a}{d^2},\frac{b}{c^2}\right].
\]
Monotonicity in each argument proves that the nested closure and
self-scaling cylinders produce nested intervals for the ratio; the
positivity of their limits and the contraction of their diameters determine
\(\pi/\varphi^2\). The scalar expression thus retains a chain of
operations compatible with depth, which is subsequently applied in
its two orientations.
